%% P02 --- Błąd numeryczny, stabilność i obliczenia w skali logarytmicznej
%%
%% Każda liczba w tym programie, której czytelnik nie policzy w pamięci,
%% pochodzi z code/p02_numerical_stability.py i trafia na stronę przez \val{}.
%% Oba listingi konsolowe to zatwierdzone transkrypty pisane przez ten skrypt.
%%
%% TEZA: algorytm może być poprawny w arytmetyce dokładnej i bezużyteczny w
%% zmiennoprzecinkowej. Program P01 powiedział, czym jest ta arytmetyka; ten
%% program zadaje pytanie inżynierskie, a odpowiedzią jest MAŁY KATALOG, a nie
%% zasada.
%%
%% CO TEN PROGRAM JEST WINIEN I SPŁACA. Siedem programów obiecuje go z nazwy i
%% każdy z nich został otwarty, zanim to powstało: F01 (ograniczenie błędu na
%% dziewięćdziesięciu sześciu warstwach), F02 (przed czym chroni epsilon w
%% normalizacji), F03 (czemu MAKSIMUM wyciąga się przed nawias i ile kosztuje
%% inny wybór; oraz logarytm prawdopodobieństwa liczony bez tworzenia
%% prawdopodobieństwa), F05 (czemu log i sigmoida to jedna operacja), F06
%% (czemu skala wierszy ma znaczenie dla solvera), F07 (dokładny wynik, przy
%% którym naiwny softmax się przepełnia) i P01 (nagromadzona strata sumowania
%% wraz z katalogiem poprawek oraz katastrofalne skracanie z nazwy).
%%
%% CZEGO F03 I F07 JUŻ NIE ZOSTAWIŁY: F03 wyprowadza dwuskładnikową tożsamość
%% log-sum-exp w całości, a F07 mówi w ramce ostrzeżenia, że odjęcie maksimum
%% jest TOŻSAMOŚCIĄ, a nie sztuczką. Oba są tu cytowane w jednym zdaniu i żadne
%% nie jest wyprowadzane ponownie. Zostaje -- i to lepsza połowa -- KTÓRY
%% element wyciągnąć i ile kosztuje zły wybór.
%%
%% CO TEN PROGRAM ŚWIADOMIE ZOSTAWIA, sprawdzone w tools/programs.json:
%%   wskaźnik uwarunkowania macierzy, źle uwarunkowane układy      -> P10
%%   FLOP-y, bajty, intensywność arytmetyczna                      -> P03
%%   czym JEST prawdopodobieństwo i czemu entropia krzyżowa         -> P26, P30
%%
%% Bliźniak: programs/en/P02-numerical-stability.tex. Ramka w ramkę, makro w
%% makro, liczba w liczbę; tools/parity.py je porównuje.

\program{Błąd numeryczny, stabilność i obliczenia w skali logarytmicznej}
\label{prog:P02}
\index{troubleshooting!an algorithm that is correct and useless}
\index{floating point!catastrophic cancellation}

\begin{outcomes}
\outcome{1--8}{powiedzieć, co czyni algorytm niebezpiecznym w arytmetyce
  zmiennoprzecinkowej, gdy jego algebrze nie można nic zarzucić}
\outcome{9--13}{powiedzieć, co odejmowanie robi z błędem względnym i czym to
  się staje na wielu warstwach}
\outcome{14--21}{policzyć wynik, przy którym eksponenta opuszcza format, i
  powiedzieć, który składnik wyciągnąć przed log-sum-exp i dlaczego}
\outcome{22--28}{powiedzieć, czemu prawdopodobieństwa nie tworzy się jako
  wartości pośredniej, i nazwać operacje, które tego unikają}
\outcome{29--37}{wybierać spośród sposobów dodawania bardzo wielu liczb i
  powiedzieć, co daje każdy z nich}
\end{outcomes}

\begin{quiz}
\item Pięć odczytów to $\val{p02.var.offset}$ i cztery kolejne mikrosekundy po
  nim. Ich wariancję liczy się w \code{float32} jako
  $\Ex[x^{2}] - (\Ex[x])^{2}$. Co wraca? \teachesat{2--3}
  \answerto{$\val{p02.var.naive}$. Wariancja nie może być ujemna, a prawdziwa
  odpowiedź to $\val{p02.var.true}$.}
\item Dwie liczby zgadzają się na dziewięciu cyfrach znaczących i zostają
  odjęte w \code{float32}. Ile cyfr znaczących ma wynik? \teachesat{5--6}
  \answerto{Żadnej. \code{float32} niesie około siedmiu, więc obie są
  przechowywane identycznie na cyfrach, na których się zgadzają, a różnica
  składa się wyłącznie z tego, co zostało po zaokrągleniu.}
\item Wynik uwagi sięga $\val{p02.over.fp16.int}$. Co robi jego eksponenta w
  \code{fp16}? \teachesat{14--15}
  \answerto{Przepełnia się do \code{inf}. Sufit zostaje przekroczony przy
  wyniku $\val{p02.over.fp16}$, a wszystko dalej staje się \code{nan}.}
\item $\ln\sum_i e^{z_i} = c + \ln\sum_i e^{z_i - c}$ zachodzi dla każdego $c$.
  Czemu zatem bierze się $c$ równe maksimum? \teachesat{17--18}
  \answerto{Bo to jedyny wybór ograniczający każdy pozostały wykładnik do zera
  lub poniżej, więc jedyny z gwarancją braku przepełnienia dla każdego wejścia,
  a nie dla tego jednego.}
\item Potrzebujesz $\ln \sigma(x)$ przy $x = \val{p02.logsig.under}$. Co daje
  policzenie najpierw $\sigma$? \teachesat{24--25}
  \answerto{\code{-inf}, bo $\sigma$ wpadła w niedomiar do
  dokładnego zera. Policzone jako jedna operacja daje
  $\val{p02.logsig.stable}$.}
\item $\val{p02.sum.n}$ wartości równych $\val{p02.sum.small}$ dodaje się po
  jednej do bieżącej sumy \code{float32}, która startuje z $1$. Ile wynosi
  suma? \teachesat{29--30}
  \answerto{$\val{p02.sum.naive32}$. Każda z nich leży poniżej połowy odstępu
  przy $1$, więc żadna nie rusza sumy.}
\item Podaj dwa sposoby naprawienia tej sumy i powiedz, który jest lepszy.
  \teachesat{31--35}
  \answerto{Dowolne dwa z: dodawanie od najmniejszych, sumowanie z kompensacją
  (Kahana), drzewo parami, szerszy akumulator. Sortowanie odzyskuje
  $\val{p02.sum.got.sorted}$\% wkładu, a pozostałe trzy ponad
  $\val{p02.sum.got.kahan}$\%.}
\item Czemu entropię krzyżową liczy się z logitów, a nie z prawdopodobieństw?
  \teachesat{22--23}
  \answerto{Bo to prawdopodobieństwo jest tą wartością pośrednią, która się
  przepełnia, wpada w niedomiar albo zaokrągla do dokładnego $0$ lub $1$;
  złożona operacja nigdy go nie tworzy i nie ma tego problemu z zakresem.}
\end{quiz}

% ============================================================
\section{Poprawny w arytmetyce dokładnej}
\label{sec:P02-correct}
\index{troubleshooting!an algorithm that is correct and useless}

\begin{fr}
Program~\ref{prog:P01} ustalił, czym jest ta arytmetyka: skończonym zbiorem
liczb, rozstawionych proporcjonalnie do swojej wielkości, z zaokrągleniem po
każdym działaniu. Ten program zadaje pytanie, które z tego wynika, i jest to
pytanie inżynierskie, a nie matematyczne. \textbf{Skoro tak, które algorytmy to
przeżyją?}

Zacznij od takiego, który nie przeżywa, a jest to algorytm, który sam pisałeś.

Pięć odczytów opóźnienia, w mikrosekundach: $\val{p02.var.offset}$ i cztery
kolejne mikrosekundy po nim. Ile wynosi ich wariancja?
\nextframe
\end{fr}

\begin{fr}
\ans{$\val{p02.var.true}$}
Cztery wartości w odległościach $-2, -1, 1, 2$ od średniej równej
$\val{p02.var.offset} + 2$, jedna w zerze, kwadraty sumujące się do $10$,
podzielone przez pięć.

Teraz obliczenie. Każdy podręcznik podaje wzór jednoprzebiegowy i jest on
naprawdę kuszący: potrzebuje jednego przejścia i dwóch bieżących sum zamiast
dwóch przejść.
\[ \Var(x) = \Ex[x^{2}] - \left(\Ex[x]\right)^{2} \]
Algebrze nie można nic zarzucić \dash{} rozwiń $(x - \mu)^2$, a wychodzi w
dwóch linijkach. Przewidz, co zwraca dla tych pięciu odczytów w \code{float32}.
\nextframe
\end{fr}

\begin{fr}
\ans{$\val{p02.var.naive}$}
\transcript{p02-variance}

\begin{trapbox}
\textbf{Wariancja jest średnią kwadratów. Nie może być ujemna, a ta wynosi
$\val{p02.var.naive}$.}

Nic w algebrze nie poszło źle. Odejmowane wielkości to
\[ \Ex[x^{2}] = \val{p02.var.ex2}, \qquad
   \left(\Ex[x]\right)^{2} = \val{p02.var.exsq} \]
i zgadzają się na $\val{p02.var.shared}$ cyfrach znaczących, podczas gdy
\code{float32} niesie ich około siedmiu. W tym formacie obie są więc
przechowywane jako liczby o identycznych cyfrach wiodących, odejmowanie skraca
każdą wspólną cyfrę, a to, co wraca, składa się z tego, co przypadkiem zostało
pod spodem \dash{} czyli z zaokrąglenia, a nie z wariancji.

Wynik to dokładnie jeden odstęp przy tej wielkości \dash{} odstęp
\code{float32} w pobliżu $\val{p02.var.ex2}$ wynosi $\val{p02.var.gap}$
\dash{} i dlatego jest to $\val{p02.var.naive}$, a nie jakaś poszarpana
liczba. To najmniejsza niezerowa odpowiedź, jaką to odejmowanie mogło dać. To
nie jest szum; to rozstaw.
\end{trapbox}
\end{fr}

\begin{fr}
Mechanizm ma nazwę i warto ją mieć, bo rozpoznasz go po kształcie na długo
przed tym, zanim będziesz umiał go dowieść: \textbf{katastrofalne skracanie}.

Odejmij dwie liczby zgodne na wiodących cyfrach znaczących, a cyfry, na których
są zgodne, skrócą się do niczego. To nie jest utrata dokładności w odejmowaniu
\dash{} odejmowanie jest dokładne i często bywa jedynym dokładnym krokiem w
całym rachunku. Strata nastąpiła wcześniej, gdy zaokrąglano każdy argument, a
odejmowanie tylko usunęło wszystko inne i ją obnażyło.

\medskip
\noindent
Reguła do zabrania: \textbf{szukaj odejmowania dwóch prawie równych wielkości i
pytaj, skąd się wzięły.}
\end{fr}


\begin{fr}
Dwie liczby zgadzają się na $\val{p02.var.shared}$ cyfrach znaczących i zostają
odjęte w \code{float32}, który niesie ich około siedmiu. Ile cyfr znaczących ma
wynik?
\nextframe
\end{fr}

\begin{fr}
\ans{Żadnej. Nie \enquote{mniej} \dash{} żadnej.}
Każda cyfra, którą trzymał \code{float32}, jest cyfrą wspólną obu argumentom,
więc każda się skraca. Cokolwiek wraca, jest zbudowane z bitów, które i tak
były już błędne.

To jest \emph{katastrofalne} z tej nazwy i jest to prawdziwe urwisko, a nie
zbocze: zgadzaj się na sześciu cyfrach w \code{float32}, a zostanie ci mniej
więcej jedna; zgadzaj się na ośmiu, a nie zostanie żadna. Liczy się wielkość
\emph{ile cyfr dzielą argumenty wobec tego, ile trzyma format}, i warto ją
policzyć, zanim uwierzysz jakiejkolwiek różnicy.
\end{fr}

\mermaidfig{p02-cancellation}{Co usuwa odejmowanie i co zostawia po
sobie.}{gdzie widać zaokrąglenie}

\begin{fr}
Lekarstwem nie jest większa precyzja. Uruchom ten sam wzór w \code{float64}, a
te odczyty przeżyje spokojnie \dash{} i polegnie na odczytach o kilka rzędów
wielkości większych, bo strukturalnie nic się nie zmieniło. \textbf{Szerszy
format przesuwa urwisko; nie usuwa go.}

Lekarstwem jest inny algorytm, a dla wariancji jest dobry. Metoda Welforda
trzyma bieżącą średnią i bieżącą sumę kwadratów odchyleń \emph{od tej
średniej}, aktualizując obie przy każdej nowej wartości:
\begin{align*}
  \mu_n &= \mu_{n-1} + \frac{x_n - \mu_{n-1}}{n} \\
  M_n &= M_{n-1} + (x_n - \mu_{n-1})(x_n - \mu_n)
\end{align*}
To wciąż jedno przejście i dwie bieżące sumy. Popatrz, czego nigdy nie tworzy,
i powiedz, czemu na tym polega cała różnica.
\nextframe
\end{fr}

\begin{fr}
\begin{ansblock}
Nigdy nie tworzy $\Ex[x^{2}]$, więc nigdy nie ma dużej wielkości, od której
trzeba odjąć dużą wielkość.
\end{ansblock}

Każde odejmowanie u Welforda to $x_n - \mu$, czyli wartość minus średnia
podobnych wartości \dash{} a to różnica rzędu rozrzutu, a nie rzędu danych.
Katastrofalne odejmowanie zostało wyprojektowane na zewnątrz, a nie
uczynione dokładniejszym, i na powyższych odczytach metoda zwraca
$\val{p02.var.welford}$ przy prawdziwej wartości $\val{p02.var.true}$.

\begin{warning}
\textbf{Taki jest kształt każdej poprawki w tym programie.} Żadna z nich nie
czyni operacji dokładniejszą. Każda usuwa operację, która niszczyła informację,
i zastępuje ją arytmetyką na wielkościach porównywalnego rzędu. Jeśli
proponowane lekarstwo tego nie robi, jest szerszym formatem w przebraniu i
polegnie później, a nie wcale.
\end{warning}
\end{fr}

% ============================================================
\section{Ile kosztuje odejmowanie i w co się kumuluje}
\label{sec:P02-cost}
\index{floating point!relative error}

\begin{fr}
Ten koszt warto podać jako wielkość, a nie jako ostrzeżenie, bo wtedy da się go
oszacować.

Załóżmy, że $a$ i $b$ niosą każde błąd względny rzędu $\eps$ \dash{} co według
reguł Programu~\ref{prog:P01} rzeczywiście robią, gdzie $\eps$ to epsilon
formatu. Ich błędy bezwzględne wynoszą wtedy około $\eps|a|$ i $\eps|b|$, a
różnica niesie ich sumę. Jako błąd \emph{względny} wyniku:
\[ \frac{\eps\left(|a| + |b|\right)}{|a - b|} \]
Przeczytaj mianownik. Gdy $a$ i $b$ są daleko od siebie, jest porównywalny z
licznikiem i nic szczególnego się nie dzieje. Gdy są blisko, dąży do zera, a
licznik nie.

Jedno założenie tego wyprowadzenia warto nazwać, bo czyni ono to wyrażenie
dolnym ograniczeniem, a nie oszacowaniem. Przyjmuje ono, że błąd względny
każdego argumentu wynosi około $\eps$, co Program~\ref{prog:P01} dopuszcza dla
liczby świeżo zaokrąglonej do formatu i dla niczego więcej. Jeśli $a$ i $b$ są
same wynikami długich obliczeń, przychodzą z większym błędem, a ten sam
współczynnik mnoży to, z czym przyszły.

Zapisz więc współczynnik wzmocnienia, a potem podaj jego rząd wielkości dla
powyższej wariancji.
\dotline
\nextframe
\end{fr}

\begin{fr}
\ans{$\frac{|a| + |b|}{|a - b|}$, a dla powyższej wariancji jest rzędu
$10^{\val{p02.var.shared}}$.}

Ten iloraz warto rozpoznawać, bo nie ogranicza się do odejmowania.

Program~\ref{prog:F06} pytał w zadaniu dodatkowym, co się dzieje, gdy jedno
równanie układu pomnożyć przez $10^{6}$: rozwiązanie się nie zmienia, para
prostych się nie zmienia, a zachowanie solvera owszem. Powód jest tą samą
wielkością. Zadanie, w którym mała względna zmiana wejścia wymusza dużą
względną zmianę wyjścia, jest \emph{źle uwarunkowane}, a uwarunkowanie jest
własnością zadania, a nie twojego kodu \dash{} żadna staranność implementacji
go nie usuwa.

\begin{rigourbox}
\textbf{Ogólną postacią tej liczby jest wskaźnik uwarunkowania}, a dla układu
równań jest on stwierdzeniem o macierzy, której ta książka jeszcze nie
zdefiniowała. Program~\ref{prog:P10} podaje go porządnie. Tutaj należy wersja
skalarna i nawyk, który ona kupuje: zanim uwierzysz policzonej różnicy, zapytaj,
jakie błędy względne miały jej wejścia i co zrobiła z nimi ta operacja.
\end{rigourbox}
\end{fr}

\begin{fr}
Program~\ref{prog:F01} zostawił temu programowi konkretny dług i uprzedził, że
odpowiedź będzie gorsza, niż się spodziewasz.

Warstwa sieci wykonuje arytmetykę, więc wyjście każdej warstwy niesie błąd
względny co najmniej $\eps$. Dziewięćdziesiąt sześć warstw stosuje to
$\val{p02.layers}$ razy. Gdyby błąd każdej warstwy pchał w tę samą stronę,
najgorszy przypadek wynosiłby $(1 + \eps)^{\val{p02.layers}} - 1$.

Policz to dla \code{fp16}, którego epsilon Program~\ref{prog:P01} podaje jako
około $10^{-3}$, i podaj odpowiedź w procentach.
\nextframe
\end{fr}

\begin{fr}
\ans{$\val{p02.grow.fp16}$\%}
\begin{center}
\begin{tabular}{lrr}
\toprule
 & najgorszy przypadek na $\val{p02.layers}$ warstwach & błędy losowego znaku \\
\midrule
\code{fp64} & $\val{p02.grow.fp64}$\% & $\val{p02.walk.fp64}$\% \\
\code{fp32} & $\val{p02.grow.fp32}$\% & $\val{p02.walk.fp32}$\% \\
\code{fp16} & $\val{p02.grow.fp16}$\% & $\val{p02.walk.fp16}$\% \\
\code{bf16} & $\val{p02.grow.bf16}$\% & $\val{p02.walk.bf16}$\% \\
\bottomrule
\end{tabular}
\end{center}

\begin{warning}
\textbf{Ograniczenie nie jest przewidywaniem, a cytowanie pierwszej kolumny
tak, jakby nim było, to błąd, dla którego ta ramka istnieje.}

Najgorszy przypadek zakłada, że każde zaokrąglenie pcha w tę samą stronę, i to
właśnie opisuje argument o kumulowaniu z Programu~\ref{prog:F01}. Zaokrąglenia
losowego znaku nie kumulują się jak liczba warstw; kumulują się jak jej
pierwiastek, i to jest druga kolumna \dash{} około $\val{p02.grow.ratio}$ razy
mniej dla \code{fp16}.

Żadna z kolumn nie jest tym, co robi twoja sieć. Pierwsza jest górnym
ograniczeniem, którego nie da się przekroczyć, a druga oszacowaniem
zakładającym niezależność, której nie sprawdziłeś. Obie ustalają kształt:
\textbf{głębokość mnoży epsilon twojego formatu przez coś pomiędzy
pierwiastkiem z głębokości a samą głębokością}, a na szesnastu bitach jest to
liczba, o którą trzeba dbać.
\end{warning}
\end{fr}

\begin{fr}
I to ustawia resztę tego programu.

Nie ma zasady, która czyni kod zmiennoprzecinkowy bezpiecznym. Jest mały
katalog sytuacji niszczących informację, a do każdej znana alternatywa, która
jej nie niszczy. Pierwszą już poznałeś: odejmowanie prawie równych wielkości,
naprawiane przez nietworzenie ich.

Pozostałe trzy to te, które spotkasz w pętli treningowej:

\begin{itemize}
\item eksponenta czegoś dużego, która opuszcza format w całości;
\item prawdopodobieństwo utworzone jako wartość pośrednia, czyli ten sam
  problem pod nazwą, której ufasz;
\item bardzo długa suma, w której małe składniki leżą pojedynczo poniżej
  progu ustalonego przez Program~\ref{prog:P01}.
\end{itemize}
\end{fr}

% ============================================================
\section{Eksponenta i jej dwa urwiska}
\label{sec:P02-cliffs}
\index{floating point!overflow and underflow}
\index{softmax!why the maximum is subtracted}

\begin{fr}
Eksponenta jest operacją, która najłatwiej wyprowadza liczbę poza jej format,
bo zamienia dodawanie wyników na mnożenie rzędów wielkości.

Program~\ref{prog:P01} policzył sufit \code{fp16} z jego budżetu wykładnika.
Wynik, przy którym $e^{z}$ przekracza ten sufit, nie jest więc nowym faktem
\dash{} jest logarytmem liczby, którą już masz.

Podaj go z dokładnością do jednego miejsca po przecinku, a potem podaj
najmniejszy całkowity wynik, który się przepełnia.
\nextframe
\end{fr}

\begin{fr}
\ans{$\val{p02.over.fp16}$, więc wynik $\val{p02.over.fp16.int}$ się przepełnia.}
\transcript{p02-cliff}

\textbf{To nie jest liczba egzotyczna.} To wynik w tej skali, w której wyniki
naturalnie występują, a utrzymanie wyników uwagi dość małymi, by były
bezpieczne, jest całym zadaniem czynnika $1/\sqrt{d_k}$, który wyprowadza
Program~\ref{prog:P25}. Awaria nie jest rzadkim zdarzeniem na skraju zakresu
formatu; jest wartością, przed którą architektura ma wbudowaną poprawkę, a
produkuje \code{inf}, a potem \code{nan}, a nie błąd, który ktokolwiek nazwałby
błędem.

I zawodzi jako skok, a nie jako zbocze, co jest przeciwieństwem tego, co
pokazała poprzednia sekcja. Kasowanie degraduje się stopniowo: współczynnik
wzmocnienia rośnie, gdy argumenty zbliżają się do siebie, więc tracisz cyfry,
potem kolejne cyfry, zanim stracisz wszystkie. Eksponenta poniżej progu jest
dokładna z dokładnością formatu, a jedną jednostkę powyżej daje \code{inf}, i
nie ma nic pomiędzy. Test, który przechodzi zakres i sprawdza błąd, przejdzie
więc czysto aż do ostatniego bezpiecznego wyniku i nie powie nic o
następnym.

Następna ramka podaje ten wynik dla wszystkich czterech formatów. Zanim to
zrobi: dwa z czterech przepełniają się przy tym samym wyniku. Powiedz które i
powiedz dlaczego.
\dotline
\nextframe
\end{fr}

\begin{fr}
\ans{\code{bf16} i \code{fp32}, bo dzielą budżet wykładnika, a zasięg
eksponenty wyznacza sam ten budżet.}

Oba urwiska we wszystkich czterech formatach \dash{} wynik, przy którym
eksponenta wychodzi górą, i wynik, przy którym wychodzi dołem:

\begin{center}
\begin{tabular}{lrr}
\toprule
 & przepełnia się powyżej & wpada w niedomiar poniżej \\
\midrule
\code{fp64} & $\val{p02.over.fp64}$ & $\val{p02.under.fp64}$ \\
\code{fp32} & $\val{p02.over.fp32}$ & $\val{p02.under.fp32}$ \\
\code{fp16} & $\val{p02.over.fp16}$ & $\val{p02.under.fp16}$ \\
\code{bf16} & $\val{p02.over.bf16}$ & $\val{p02.under.bf16}$ \\
\bottomrule
\end{tabular}
\end{center}

To powtórzenie pointy Programu~\ref{prog:P01}, a tabela mówi też resztę:
\textbf{mantysa nie pojawia się w niej nigdzie.} Żadna z kolumn nie drgnie,
gdy mantysa się zmieni, po żadnej ze stron.
\end{fr}

\mermaidfig{p02-two-cliffs}{Dwa sposoby, na jakie liczba opuszcza format, i co
ustala odległość między nimi.}{dwa urwiska}

\begin{fr}
Program~\ref{prog:F07} rozstrzygnął już, co robić z górnym urwiskiem, i
rozstrzygnął to algebrą, a nie radą: odjęcie stałej od każdego wyniku przed
eksponentą nie zmienia softmaksu, bo eksponenta tej stałej skraca się w
liczniku i mianowniku jednakowo. Jego ramka ostrzeżenia mówi wprost \dash{}
\textbf{to jest tożsamość, a wersja, która tego nie robi, jest tą błędną.}

Program~\ref{prog:F03} zrobił ten sam krok dla pary logarytmów, wyciągając
większy z nich przed $\ln(e^{u} + e^{v})$.

Oba zostawiły to samo pytanie otwarte i to ono jest warte sekcji. Ogólnie
zapisana tożsamość brzmi
\[ \ln \sum_i e^{z_i} = c + \ln \sum_i e^{z_i - c} \]
i \emph{zachodzi dla każdego $c$}. Nic powyżej nie wymusza wyboru. Czemu więc
maksimum jest tym, którego używa każda biblioteka?
\nextframe
\end{fr}

\begin{fr}
\begin{ansblock}
Bo jest jedynym wyborem ograniczającym każdy pozostały wykładnik do zera lub
poniżej, a więc jedynym z gwarancją braku przepełnienia \dash{} dla każdego
wejścia, a nie dla tego jednego.
\end{ansblock}

Przy $c = \max_i z_i$ każde $z_i - c$ jest co najwyżej $0$, więc każde
$e^{z_i - c}$ jest co najwyżej $1$ i żaden składnik nie przekroczy sufitu
żadnego formatu. Suma co najwyżej $n$ składników, każdego co najwyżej $1$,
wynosi co najwyżej $n$, z czym też żaden format nie ma kłopotu.

Weź dowolne inne $c$, a największy tworzony składnik to $e^{\max_i z_i - c}$,
które jest powyżej $1$ i rośnie wykładniczo w miarę spadku $c$. \textbf{To, czy
się mieści, jest wtedy własnością twoich danych.}
\end{fr}

\mermaidfig{p02-pivot}{Przed nawias można wyciągnąć dowolną stałą; zmienia się
największy składnik, który trzeba utworzyć.}{wybór elementu}

\begin{fr}
To ostatnie zdanie warte jest pomiaru, bo wyjaśnia, czemu ten wybór tak łatwo
uchodzi na sucho.

Weź wiersz $\val{p02.pivot.total}$ wyników i wyciągaj kolejno każdy z nich, w
\code{fp16}. Ile z $\val{p02.pivot.total}$ możliwości utrzymuje każdy składnik
wewnątrz formatu?
\nextframe
\end{fr}

\begin{fr}
\ans{$\val{p02.pivot.safe}$ z nich.}
\begin{trapbox}
\textbf{Większość błędnych wyborów działa.} \code{fp16} znosi niedobór do
$\val{p02.over.fp16}$ między największym wynikiem a wyciąganym, a w zwyczajnym
wierszu większość wpisów mieści się w tym zakresie. Naiwna implementacja, która
bierze pierwszy wynik albo ostatni, albo nie wyciąga niczego w wierszu, którego
wyniki akurat są małe, działa poprawnie \dash{} przez zestaw testów, przez
przegląd kodu i dalej na produkcję.

Najgorszy wybór w tym wierszu musi utworzyć składnik równy
$\val{p02.pivot.worst}$, co nie jest blisko. Ale jest to najgorszy wybór
\emph{w tym wierszu}: zmień dane, a granica się przesunie, bo zawsze była
faktem o danych.

\textbf{Maksimum jest jedynym wyborem, który nim nie jest.} To właśnie znaczy
\enquote{stabilny numerycznie} jako termin techniczny \dash{} nie
\emph{dokładniejszy}, lecz \emph{bezpieczny dla każdego wejścia, a nie dla
tych, które akurat sprawdziłeś}.
\end{trapbox}
\end{fr}

\begin{fr}
Dolne urwisko potrzebuje jednej ramki, bo poprawka na górne załatwia je przy
okazji.

Po odjęciu maksimum najmniejsze składniki to $e^{z_i - \max}$ dla wyników
najdalszych od szczytu, a te wpadają w niedomiar do zera, gdy wynik spadnie
więcej niż $\val{p02.drop.fp16}$ poniżej maksimum w \code{fp16}. To nie jest
awaria: taki składnik i tak wniósłby mniej niż $\val{p02.drop.fp16.frac}$
największego, a jego pominięcie nie zmienia niczego mierzalnego.

\textbf{Niedomiar w sumie jest zwykle nieszkodliwy, a przepełnienie nigdy}, i
warto to wiedzieć, bo mówi, którą stronę chronić. Składnik, który staje się
zerem, to składnik, którego nie potrzebowałeś. Składnik, który staje się
\code{inf}, zabiera ze sobą sumę, logarytm, stratę i każdy gradient.

Cała ta operacja \dash{} odjąć maksimum, wziąć eksponentę, zsumować, wziąć
logarytm, dodać maksimum z powrotem \dash{} to $\logsumexp$, a dla powyższego
wiersza daje $\val{p02.lse.exact}$.
\end{fr}

% ============================================================
\section{Nigdy nie twórz prawdopodobieństwa}
\label{sec:P02-logits}
\index{softmax!computed from logits}
\index{loss!cross-entropy from logits}

\begin{fr}
Program~\ref{prog:F03} obiecał, że ten program pokaże, jak logarytm
prawdopodobieństwa tokenu liczy się \emph{bez tworzenia prawdopodobieństwa}, i
to sformułowanie było celowe. To ta sama instrukcja co w poprzedniej sekcji, o
poziom wyżej.

Prawdopodobieństwo to liczba między $0$ a $1$. Program~\ref{prog:P01} mówi, ile
ten zakres kosztuje: poniżej $\val{p02.under.fp16}$ w wykładniku wartość
\code{fp16} jest zerem, a prawdopodobieństwo blisko $1$ ma swoje interesujące
cyfry dokładnie tam, gdzie liczba zmiennoprzecinkowa ma ich najmniej, bo odstęp
przy $1$ jest największym odstępem poniżej $1$.

Prawdopodobieństwo jest więc złą wartością pośrednią na obu końcach, a jego
logarytm dobrą na obu. \textbf{Reguła brzmi: prawdopodobieństwo nie powinno
istnieć.}

Dla softmaksu oba kroki składają się w jedną operację, a algebra, która to
robi, ma trzy linijki. Rozpisz $\ln \softmax(z)_j$ i uprość.
\dotline
\nextframe
\end{fr}

\begin{fr}
\ans{$z_j - \logsumexp(z)$: żadnej eksponenty surowego wyniku, żadnego
dzielenia i nigdzie żadnego prawdopodobieństwa.}

\[ \ln \softmax(z)_j = \ln \frac{e^{z_j}}{\sum_k e^{z_k}}
   = z_j - \ln \sum_k e^{z_k} = z_j - \logsumexp(z) \]
Log-softmax to odjęcie wielkości z poprzedniej sekcji od wyniku, a odejmowanie
odbywa się między liczbami tego samego rzędu, czyli tak,
jak kazała szukać \S\ref{sec:P02-correct}.

To właśnie ma na myśli biblioteka, gdy mówi o funkcji przyjmującej
\emph{logity} zamiast prawdopodobieństw, i dlatego podanie jej
prawdopodobieństw \dash{} przez uczynne zastosowanie softmaksu wcześniej
\dash{} jest usterką, a nie nadmiarowością.
\end{fr}

\begin{fr}
Przypadek dwuklasowy ma ten sam kształt i własną nazwę, a
Program~\ref{prog:F05} na niego wskazał: $\ln \sigma(x)$, czyli log-sigmoida,
która pojawia się w każdej dwuklasowej entropii krzyżowej, której wejściem jest
logit.

Program~\ref{prog:F05} powiedział, że złożenie jest bezpieczne matematycznie i
niebezpieczne zmiennoprzecinkowo. Powiedz, mniej więcej przy jakim $x$ postać
złożona zawodzi i co wtedy zwraca.
\nextframe
\end{fr}

\begin{fr}
\ans{Poniżej mniej więcej $\val{p02.logsig.under}$ zwraca \code{-inf}.}
$\sigma(x)$ zachowuje się jak $e^{x}$ dla dużych ujemnych $x$, więc wpada w
niedomiar do dokładnego zera, gdy $x$ przekroczy $\val{p02.under.fp64}$
\dash{} to dolne urwisko z poprzedniej sekcji, przychodzące pod inną nazwą. A
$\ln 0$ to \code{-inf}, które przenosi się do straty, a potem do
każdego gradientu.

Policzone jako jedna operacja nie sprawia najmniejszego kłopotu:
\[ \ln \sigma(x) = -\ln\!\left(1 + e^{-x}\right) \]
co dla ujemnych $x$ przekształca się do $x - \ln(1 + e^{x})$, żeby argument
eksponenty nigdy nie był dodatni. Przy $x = \val{p02.logsig.under}$ zwraca
$\val{p02.logsig.stable}$, co jest poprawną odpowiedzią i ledwie odróżnialne od
samego $x$.

To był dolny koniec. Górny też zawodzi i to jest ta połowa, którą się pomija:
dwuklasowa entropia krzyżowa potrzebuje $\ln p$ i $\ln(1 - p)$, a
Program~\ref{prog:F07} zmierzył, gdzie $\sigma$ zaokrągla się do dokładnego
$\num{1.0}$. Powiedz, co robi tam strata.
\dotline
\nextframe
\end{fr}

\begin{fr}
\ans{$1 - p$ jest dokładnie zerem, więc $\ln(1 - p)$ to \code{-inf} \dash{}
przy predykcji pewnej i \emph{trafnej}.}

\begin{trapbox}
\textbf{I to jest ta połowa, którą się pomija.}

Program~\ref{prog:F07} zmierzył, że $\sigma$ zaokrągla się do dokładnego
$\num{1.0}$ mniej więcej od $\val{f07.sig.saturates}$ w górę; pierwszą liczbą
całkowitą, przy której to robi, jest $\val{p02.logsig.saturates}$. Dwuklasowa
entropia krzyżowa potrzebuje i $\ln p$, i $\ln(1 - p)$, a w tym punkcie
$1 - p$ jest dokładnym zerem, więc $\ln(1 - p)$ wynosi
\code{-inf}.

Pewna \emph{poprawna} predykcja daje więc to samo \code{-inf} co
pewna błędna, a strata staje się \code{nan} na przykładzie, który model
odgadł dobrze. To jest ta awaria, która wysyła ludzi na poszukiwanie błędu w
danych.

\textbf{Oba końce usuwa ten sam ruch}: zachowaj logit i pozwól funkcji straty
wykonać złożenie. Biblioteka dostarcza jedną funkcję właśnie dlatego, że nie ma
bezpiecznego sposobu zapisania jej jako dwóch.
\end{trapbox}
\end{fr}

\begin{fr}
\begin{aibox}
Konsekwencja w pętli treningowej, podana jako reguła, którą sprawdzisz czytając
diffa.

Przejście w przód klasyfikatora kończy się na \emph{logitach} i nic za tym
punktem nie tworzy prawdopodobieństwa, dopóki coś poza stratą nie potrzebuje go
wyświetlić. Strata dostaje logity i cele; metryka może dostać
prawdopodobieństwa, bo metryka mylona o zaokrąglenie jest usterką kosmetyczną,
a strata równa \code{nan} jest zatrzymanym przebiegiem.

Sygnałem w kodzie jest softmax albo sigmoida bezpośrednio przed logarytmem albo
przed stratą, której nazwa wspomina logity. Oba to dwuoperacyjna postać czegoś,
co biblioteka dostarcza jako jedno, i oba działają na wsadach, które
sprawdziłeś.
\end{aibox}
\end{fr}

\begin{fr}
Jedna granica, żeby ta sekcja nie została odczytana jako coś więcej, niż jest.

Nic tutaj nie mówi, czym prawdopodobieństwo \emph{jest}, czemu jego logarytm
jest właściwą rzeczą do sumowania ani czemu entropia krzyżowa jest stratą, a
nie wygodnym wyborem. Program~\ref{prog:P26} wyprowadza tę stratę z metody
największej wiarygodności, a Program~\ref{prog:P30} nadaje entropii krzyżowej
znaczenie w bitach.

Do tego programu należy coś węższego i jest to fakt o arytmetyce: \textbf{z
tych dwóch wielkości to logarytm jest tą, którą twoja maszyna pomieści.} To
prawda niezależnie od tego, co ta strata się okaże znaczyć.
\end{fr}

% ============================================================
\section{Dodawanie bardzo wielu liczb}
\label{sec:P02-summation}
\index{floating point!summation order}

\begin{fr}
Program~\ref{prog:P01} ustalił próg i przekazał konsekwencję tutaj, z nazwy. To
jest ta sekcja.

Wkład poniżej połowy odstępu przy bieżącej sumie rusza ją o dokładnie nic. Weź
więc $\val{p02.sum.n}$ wartości równych $\val{p02.sum.small}$ i dodaj je po
jednej do sumy \code{float32}, która startuje z $1$. Dokładna odpowiedź to
$\val{p02.sum.exact}$.

Co zwróci ta pętla?
\nextframe
\end{fr}

\begin{fr}
\ans{$\val{p02.sum.naive32}$. Przepadła każda z nich, co do jednej.}
Nie większość i nie na skutek zaokrąglenia: odstęp przy $1$ w \code{float32}
wynosi około $10^{-7}$, jego połowa to około $\val{p01.swamp.fp32}$, a
$\val{p02.sum.small}$ leży poniżej. Każde dodawanie jest pustą operacją, która
kończy się powodzeniem, milion razy z rzędu.

\begin{warning}
\textbf{Nie ma błędu, nie ma ostrzeżenia i nie ma \code{nan}.} Pętla się
wykonuje, suma jest zupełnie zwyczajną liczbą, a odpowiedź jest błędna o
$\val{p02.sum.err.naive32}$ \dash{} czyli o całość tego, co sumowałeś.
Akumulacja gradientu, która to robi, raportuje zdrową stratę i niczego się nie
uczy, a objawem jest model, który nie chce się trenować, a nie program, który
nie chce się wykonać.
\end{warning}
\end{fr}

\begin{fr}
Katalog ma cztery pozycje i nie są one równoważne. Pierwsza jest darmowa.

\textbf{Dodaj małe wartości do siebie nawzajem, zanim dodasz je do dużej.}
Sortowanie rosnąco to załatwia: milion drobnych wartości gromadzi się między
sobą, gdzie są porównywalnego rzędu, a suma spotyka $1$ dopiero na końcu, gdy
jest już dość duża, by to spotkanie przeżyć.

To brak łączności z trzech wartości Programu~\ref{prog:P01} w skali,
uruchomiony celowo, a nie przypadkiem. Ile spodziewasz się odzyskać?
\nextframe
\end{fr}

\begin{fr}
\ans{$\val{p02.sum.got.sorted}$\% wkładu \dash{} prawie całość, ale nie całość.}
Niedobór wynosi $\val{p02.sum.err.sorted}$ i nie jest jednym zaokrągleniem:
jest milionem zaokrągleń. Każde dodawanie do rosnącej sumy zaokrągla, błędy w
większości mają wspólny znak, a przy $\val{p02.sum.n}$ z nich ten dryf staje
się widoczny.

\textbf{Sortowanie usuwa katastrofę i zostawia dryf}, co jest uczciwym opisem i
warto go mieć, bo sortowanie to poprawka, po którą ludzie sięgają i na której
poprzestają.

Dryf bierze się stąd, że każde dodawanie zaokrągla, a zaokrąglenie jest
wyrzucane. Co trzeba by trzymać obok sumy, żeby przestać je wyrzucać?
\dotline
\nextframe
\end{fr}

\begin{fr}
\ans{Tę część ostatniego dodawania, która przepadła: jedną dodatkową zmienną,
dodawaną z powrotem do następnego składnika przed jego użyciem.}

Druga pozycja atakuje dryf wprost. \emph{Sumowanie z kompensacją} \dash{}
Kahana \dash{} trzyma dokładnie taką zmienną:
\begin{align*}
  y &= x_i - c \\
  t &= s + y \\
  c &= (t - s) - y \\
  s &= t
\end{align*}
Przeczytaj trzecią linijkę. $t - s$ to tyle, o ile dodawanie \emph{naprawdę}
ruszyło sumę, a $y$ to tyle, o ile miało ją ruszyć; różnica jest dokładnie tym,
co zostało zaokrąglone, i odzyskuje się ją odejmowaniem dwóch prawie równych
liczb.

\textbf{Skracanie jest tu mechanizmem, użytym celowo.} Ta sama operacja, która
zniszczyła wariancję w \S\ref{sec:P02-correct}, jest tutaj jedyną, która w
ogóle widzi utracone bity.

Trzecia pozycja nie potrzebuje żadnej dodatkowej zmiennej. Dodawaj w
zrównoważonym drzewie, a nie w linii: sparuj wartości, dodaj pary, sparuj te
wyniki i tak dalej. Liczba dodawań się nie zmienia. Powiedz, co się zmienia i
o ile, dla miliona wartości.
\dotline
\nextframe
\end{fr}

\begin{fr}
\ans{\emph{Głębokość}. Każda wartość dociera teraz do sumy przez mniej więcej
dwadzieścia dodawań, a nie przez milion.}

\textbf{Dodawaj w zrównoważonym drzewie}, a nie w linii, a dryf ma dwadzieścia
okazji do nagromadzenia wzdłuż dowolnej ścieżki zamiast miliona.

Sumowanie parami to właśnie to, co robi redukcja biblioteki tablicowej, i
dlatego naiwna pętla Pythona po tablicy bywa \emph{mniej} dokładna niż wywołanie
biblioteki, które zastąpiła, a nie tylko wolniejsza. Która biblioteka i która
wersja to fakt o wydaniu, a nie o matematyce, więc ta książka żadnej nie
wymienia.

Czwarta nie wymaga żadnego sprytu: \textbf{akumuluj w formacie szerszym niż
wartości.} Wartości szesnastobitowe sumowane do akumulatora \code{float32} albo
wartości \code{float32} do \code{float64}. Kosztuje odrobinę pamięci i zero
myślenia, i tak właśnie robi trening w mieszanej precyzji.

Cztery poprawki, a następna ramka ocenia je wszystkie na tym samym milionie
wartości. Zanim to zrobi: która z czterech wypadnie twoim zdaniem najgorzej?
\dotline
\nextframe
\end{fr}

\begin{fr}
\ans{Sortowanie \dash{} najtańsza z czterech i ta, którą najczęściej opisuje
się jako wystarczającą.}

Wszystkie pięć, na tym samym milionie wartości:

\begin{center}
\begin{tabular}{lrr}
\toprule
 & suma & z wkładu \\
\midrule
naiwnie \code{float32} & $\val{p02.sum.naive32}$ & $\val{p02.sum.got.naive32}$\% \\
od najmniejszych & $\val{p02.sum.sorted}$ & $\val{p02.sum.got.sorted}$\% \\
z kompensacją & $\val{p02.sum.kahan}$ & $\val{p02.sum.got.kahan}$\% \\
parami & $\val{p02.sum.pairwise}$ & $\val{p02.sum.got.pairwise}$\% \\
szerszy akumulator & $\val{p02.sum.wide}$ & $\val{p02.sum.got.wide}$\% \\
\bottomrule
\end{tabular}
\end{center}

Dokładna wartość to $\val{p02.sum.exact}$, policzona na liczbach całkowitych,
żeby to, do czego porównuje się każdy wiersz, samo nie było wynikiem
zmiennoprzecinkowym.

Dwie rzeczy w tej tabeli są warte więcej niż sam ranking. Pierwsza to przepaść
między pierwszym wierszem a każdym innym, czyli różnica między odpowiedzią a
brakiem odpowiedzi. Druga to rozmiar niedoboru, który przewidziałeś:
\textbf{sortowanie wypada o rząd wielkości gorzej niż pozostałe trzy}, więc
najtańsza poprawka nie jest po prostu najsłabszą z czterech \dash{} nie jest w
tej samej klasie co pozostałe trzy.
\end{fr}

\begin{fr}
\begin{aibox}
Program~\ref{prog:F02} zostawił tu pytanie: przed czym chroni $\eps$ w
normalizacji warstwowej i ile z tego dany format jest w stanie wchłonąć.

Chroni przed dzieleniem przez odchylenie standardowe, które wyszło zerem
\dash{} co zdarza się, gdy cecha jest stała w całym wsadzie, a także gdy jest
prawie stała, a wariancję policzono wzorem, od którego ten program się zaczął.
$\eps$ czyni mianownik dodatnim, cokolwiek zrobił licznik.

Ile format wchłania, jest teraz liczbą, którą umiesz podać: $\eps$ musi być
większe niż własna podłoga szumu formatu w skali tej wariancji, więc $\eps$
równe $10^{-8}$ nie robi zupełnie nic w \code{fp16}, którego epsilon jest rzędu
$10^{-3}$. Dlatego przepisy na mieszaną precyzję liczą statystyki normalizacji
w \code{float32}, nawet gdy wszystko wokół ma szesnaście bitów.
\end{aibox}
\end{fr}

\begin{fr}
Ostatnia ramka, i jest listą kontrolną, a nie konkluzją.

Cztery sytuacje niszczą informację, a każda ma znaną alternatywę:

\begin{itemize}
\item \textbf{odejmowanie prawie równych wielkości} \dash{} znajdź, skąd się
  wzięły, i policz różnicę wprost, jeśli możesz;
\item \textbf{eksponenta surowego wyniku} \dash{} odejmij maksimum, zawsze, i
  bierz logarytm, a nie iloraz;
\item \textbf{prawdopodobieństwo jako wartość pośrednia} \dash{} zachowaj logit
  i pozwól stracie wykonać złożenie;
\item \textbf{długa suma} \dash{} sparuj ją, skompensuj albo poszerz
  akumulator, i nie poprzestawaj na sortowaniu.
\end{itemize}

Żadna z tych czterech nie jest o byciu starannym. Każda to konkretna operacja z
konkretnym zamiennikiem, a umiejętnością jest rozpoznanie operacji, a nie
zapamiętanie poprawki. Program~\ref{prog:P03} zadaje o tej samej arytmetyce
następne pytanie \dash{} nie czy jest poprawna, ale ile kosztuje.
\end{fr}

% ============================================================
\begin{summarybox}
\sumitem{1--3}{\result{\textbf{Algorytm może być poprawny w arytmetyce dokładnej i
  bezużyteczny w zmiennoprzecinkowej.} Jednoprzebiegowa wariancja
  $\Ex[x^{2}] - (\Ex[x])^{2}$ na pięciu odczytach przy $\val{p02.var.offset}$
  zwraca $\val{p02.var.naive}$ w \code{float32}, przy prawdziwej wartości
  $\val{p02.var.true}$. Algebrze nie można nic zarzucić, a wynik nie jest
  wariancją}.}
\sumitem{4--6}{\result{\textbf{Katastrofalne skracanie} jest mechanizmem: obie
  wielkości zgadzają się na $\val{p02.var.shared}$ cyfrach znaczących, podczas
  gdy format trzyma siedem, więc odejmowanie skraca każdą wspólną cyfrę i
  oddaje zaokrąglenie. Odejmowanie jest dokładne; strata nastąpiła przy
  zaokrąglaniu argumentów. Szukaj różnicy prawie równych wielkości i pytaj,
  skąd się wzięły}.}
\sumitem{7--8}{\result{\textbf{Lekarstwem nigdy nie jest większa precyzja} \dash{}
  szerszy format przesuwa urwisko, zamiast je usuwać. Metoda Welforda nigdy nie
  tworzy $\Ex[x^{2}]$, więc każde odejmowanie jest rzędu rozrzutu, a nie rzędu
  danych, i zwraca $\val{p02.var.welford}$. Każda poprawka w tym programie ma
  ten kształt: usuń niszczącą operację, nie czyń jej dokładniejszą}.}
\sumitem{9--10}{\result{Wzmocnienie wynosi $\frac{|a| + |b|}{|a - b|}$ i jest
  własnością \emph{zadania}. Jego ogólną postacią jest wskaźnik uwarunkowania,
  który należy do Programu~\ref{prog:P10}, bo jest stwierdzeniem o macierzy}.}
\sumitem{11--13}{\result{Na $\val{p02.layers}$ warstwach najgorszy przypadek
  \code{fp16} to $\val{p02.grow.fp16}$\%, a \code{bf16} to
  $\val{p02.grow.bf16}$\% \dash{} ale \textbf{ograniczenie nie jest
  przewidywaniem}: błędy losowego znaku rosną jak pierwiastek z głębokości,
  około $\val{p02.grow.ratio}$ razy mniej. Głębokość mnoży epsilon przez coś
  pomiędzy $\sqrt{n}$ a $n$}.}
\sumitem{14--16}{\result{\textbf{Eksponenta opuszcza format przy logarytmie jego
  sufitu}, czyli przy $\val{p02.over.fp16}$ dla \code{fp16} \dash{} więc wynik
  $\val{p02.over.fp16.int}$ się przepełnia, a to zwyczajny logit uwagi.
  \code{bf16} i \code{fp32} dzielą ten wiersz przy $\val{p02.over.fp32}$, bo
  zakres to sam budżet wykładnika}.}
\sumitem{17--20}{\result{$\ln\sum_i e^{z_i} = c + \ln\sum_i e^{z_i - c}$ zachodzi dla
  \emph{każdego} $c$, więc tożsamość nie wymusza wyboru. \textbf{Maksimum jest
  jedynym $c$ ograniczającym każdy wykładnik do zera lub poniżej}, a więc
  jedynym bezpiecznym dla każdego wejścia. Zmierzone: $\val{p02.pivot.safe}$ z
  $\val{p02.pivot.total}$ wyborów działa w jednym zwyczajnym wierszu, i dlatego
  zły wybór przechodzi przez przegląd \dash{} a najgorszy tworzy
  $\val{p02.pivot.worst}$}.}
\sumitem{21}{\result{\textbf{Niedomiar w sumie jest zwykle nieszkodliwy, a
  przepełnienie nigdy.} Składnik, który staje się zerem, i tak nic nie wnosił;
  składnik, który staje się \code{inf}, zabiera stratę i każdy gradient. Cała
  ta operacja to $\logsumexp$}.}
\sumitem{22--23}{\result{\textbf{Prawdopodobieństwo jest złą wartością pośrednią na obu
  końcach}, a jego logarytm dobrą, więc prawdopodobieństwo nie powinno
  istnieć. $\ln\softmax(z)_j = z_j - \logsumexp(z)$: żadnej eksponenty surowego
  wyniku, żadnego dzielenia i odejmowanie między liczbami tego samego rzędu}.}
\sumitem{24--26}{\result{$\ln\sigma(x)$ policzone w dwóch krokach zwraca
  \code{-inf} poniżej mniej więcej $\val{p02.logsig.under}$,
  gdzie $\sigma$ wpadła w niedomiar; jako jedna operacja zwraca
  $\val{p02.logsig.stable}$. \textbf{Górny koniec też zawodzi}: $\sigma$ jest
  dokładnym $\num{1.0}$ od $\val{p02.logsig.saturates}$, więc $\ln(1-p)$ to
  \code{-inf}, a pewna \emph{poprawna} predykcja produkuje
  \code{nan}}.}
\sumitem{27--28}{\result{Przejście w przód kończy się na logitach, a strata wykonuje
  złożenie. Sygnałem w diffie jest softmax albo sigmoida bezpośrednio przed
  logarytmem. Czym prawdopodobieństwo \emph{jest} i czemu entropia krzyżowa
  jest stratą, należy do Programów~\ref{prog:P26} i~\ref{prog:P30}; tutaj
  należy to, że z tych dwóch wielkości logarytm jest tą, którą maszyna
  pomieści}.}
\sumitem{29--30}{\result{$\val{p02.sum.n}$ wartości równych $\val{p02.sum.small}$
  dodanych do sumy \code{float32} równej $1$ daje $\val{p02.sum.naive32}$
  \dash{} \textbf{przepadła każda}, bez błędu i bez ostrzeżenia, z pomyłką o
  całość sumowanej wielkości}.}
\sumitem{31--35}{\result{Cztery poprawki, i nie są równoważne: od najmniejszych
  odzyskuje $\val{p02.sum.got.sorted}$\%, kompensacja i drzewo parami
  $\val{p02.sum.got.kahan}$\%, szerszy akumulator
  $\val{p02.sum.got.wide}$\%. \textbf{Sortowanie jest najtańsze i najsłabsze},
  a to je najczęściej nazywa się wystarczającym. Trzecia linijka Kahana
  odzyskuje utracone bity skracaniem użytym celowo}.}
\sumitem{36--37}{\result{$\eps$ w normalizacji warstwowej chroni dzielenie przez
  odchylenie standardowe, które wyszło zerem, i musi przekraczać podłogę szumu
  formatu w skali tej wariancji \dash{} dlatego statystyki normalizacji liczy
  się w \code{float32} wewnątrz modelu szesnastobitowego. \textbf{Umiejętnością
  jest rozpoznanie operacji, a nie zapamiętanie poprawki.}}}
\end{summarybox}

\canyou

% ============================================================
\begin{testexercises}
\item Powiedz, czym jest katastrofalne skracanie i w którym miejscu rachunku
  informacja naprawdę przepadła.
  \answerto{Odejmowaniem dwóch liczb zgodnych na wiodących cyfrach znaczących,
  więc te cyfry się skracają, a zostaje to, co zbudowały bity niskiego rzędu.
  Samo odejmowanie jest dokładne; strata nastąpiła przy zaokrąglaniu każdego
  argumentu.}
\item Wariancja jednoprzebiegowa zwróciła $\val{p02.var.naive}$ przy prawdziwej
  wartości $\val{p02.var.true}$. Powiedz, czemu uruchomienie jej w
  \code{float64} nie jest lekarstwem.
  \answerto{Przesuwa przesunięcie, przy którym wzór zawodzi, i nie zmienia
  niczego strukturalnie. Odejmowanie dwóch prawie równych dużych wielkości
  wciąż tam jest, więc algorytm wciąż niszczy informację dla dostatecznie
  dużych danych.}
\item Podaj współczynnik wzmocnienia przy liczeniu $a - b$ z argumentów
  niosących każdy błąd względny $\eps$ i powiedz, od czego zależy.
  \answerto{$\frac{|a| + |b|}{|a - b|}$, razy $\eps$. Zależy od tego, jak
  blisko siebie są argumenty, co jest własnością zadania, a nie kodu.}
\item Podaj wynik, przy którym eksponenta przepełnia \code{fp16}, i wynik, przy
  którym przepełnia \code{bf16}, i powiedz, czemu tak bardzo się różnią.
  \answerto{$\val{p02.over.fp16}$ i $\val{p02.over.bf16}$. Zakres eksponenty to
  logarytm sufitu formatu, a ten ustala sam budżet wykładnika \dash{} pięć
  bitów przeciw ośmiu.}
\item $\ln \sum_i e^{z_i} = c + \ln \sum_i e^{z_i - c}$. Powiedz, co ogranicza
  $c$ matematycznie, a co w praktyce.
  \answerto{Matematycznie nic \dash{} tożsamość zachodzi dla każdego $c$. W
  praktyce tylko $c = \max_i z_i$ ogranicza każdy pozostały wykładnik do zera
  lub poniżej, więc tylko ono jest bezpieczne dla każdego wejścia.}
\item Kolega mówi, że jego softmax wyciąga pierwszy wynik zamiast maksimum i
  nigdy nie zawiódł. Powiedz, co z tego wynika.
  \answerto{Że jego wiersze nigdy nie były rozrzucone bardziej niż o około
  $\val{p02.over.fp16}$ w \code{fp16}. To fakt o danych, które widział, a nie o
  kodzie, a granica przesuwa się w chwili, gdy przesuwają się dane.}
\item Powiedz, czemu podanie prawdopodobieństw funkcji straty oczekującej
  logitów jest usterką, a nie równoważną drogą.
  \answerto{To prawdopodobieństwo jest tą wartością pośrednią, która ma problem
  z zakresem: nie przepełnia się, ale wpada w niedomiar do zera i zaokrągla do
  dokładnej jedynki. Postać przyjmująca logity nigdy go nie tworzy i składa oba
  kroki w arytmetykę na liczbach porównywalnego rzędu.}
\item Klasyfikator dwuklasowy daje stratę \code{nan} na przykładzie, którego
  jest pewny i który odgadł dobrze. Podaj mechanizm.
  \answerto{$\sigma$ zaokrągliła się do dokładnego $\num{1.0}$, co dzieje się
  od $\val{p02.logsig.saturates}$ w górę, więc $1 - p$ jest dokładnym zerem, a
  $\ln(1-p)$ to \code{-inf}. Liczenie straty z logitu to usuwa.}
\item $\val{p02.sum.n}$ wartości równych $\val{p02.sum.small}$ dodaje się do
  sumy \code{float32} równej $1$. Podaj wynik i powiedz, jakie dostaniesz
  ostrzeżenie.
  \answerto{$\val{p02.sum.naive32}$, i żadnego. Każda wartość leży poniżej
  połowy odstępu przy $1$, więc każde dodawanie kończy się powodzeniem i nic
  nie rusza.}
\item Uszereguj cztery poprawki sumowania według tego, ile odzyskują, i
  powiedz, którą najczęściej się przecenia.
  \answerto{Szerszy akumulator, potem kompensacja i drzewo parami razem, potem
  od najmniejszych przy $\val{p02.sum.got.sorted}$\%. Sortowanie jest tą
  opisywaną jako wystarczająca, a wypada o rząd wielkości gorzej niż pozostałe
  trzy.}
\end{testexercises}

% ============================================================
\begin{furtherproblems}
\item Wyprowadź $\Var(x) = \Ex[x^{2}] - (\Ex[x])^{2}$ z definicji i powiedz, w
  której linijce wprowadzasz kłopot numeryczny.
  \answerto{Rozwiń $\Ex[(x - \mu)^{2}] = \Ex[x^{2}] - 2\mu\Ex[x] + \mu^{2}$ i
  użyj $\Ex[x] = \mu$. Kłopot przychodzi wraz z samym rozwinięciem: zastępuje
  sumę małych kwadratów odchyleń różnicą dwóch dużych wielkości, a te są równe
  w arytmetyce dokładnej i nie są równe na maszynie.}
\item Pokaż, że aktualizacja Welforda utrzymuje
  $M_n = \sum_{i \le n} (x_i - \mu_n)^2$, i powiedz, czemu dwie różne średnie w
  jej drugiej linijce nie są literówką.
  \answerto{Do sprawdzenia jest tożsamość
  $M_n - M_{n-1} = (x_n - \mu_{n-1})(x_n - \mu_n)$, która wynika z
  $\mu_n - \mu_{n-1} = (x_n - \mu_{n-1})/n$. Stara i nowa średnia pojawiają się
  dlatego, że poprawka musi uwzględnić przesunięcie każdego wcześniejszego
  odchylenia, gdy średnia się rusza, a ich iloczyn robi dokładnie to.}
\item Wiersz wyników jest rozrzucony w zakresie $200$. Powiedz, co się dzieje
  pod softmaksem odejmującym maksimum, w \code{fp16}.
  \answerto{Górne składniki są w porządku, a wszystko ponad
  $\val{p02.drop.fp16}$ poniżej maksimum wpada w niedomiar do zera, więc te
  tokeny dostają prawdopodobieństwo dokładnie zerowe. Dostawały mniej niż
  $\val{p02.drop.fp16.frac}$ największego składnika, więc nic mierzalnego nie
  ginie \dash{} ale ich gradienty są dokładnie zerowe, a nie bardzo małe.}
\item Wyjaśnij, czemu $\logsumexp$ wiersza $n$ równych wyników to ten wynik
  plus $\ln n$, i użyj tego do sprawdzenia implementacji.
  \answerto{Odjęcie maksimum zostawia $n$ składników równych $1$, których suma
  wynosi $n$; dodanie maksimum z powrotem daje $z + \ln n$. To najtańszy
  możliwy test i natychmiast wykrywa błąd wyciągania, bo zły wybór wciąż daje
  tu poprawną odpowiedź.}
\item Entropię krzyżową liczy się jako \code{-log(softmax(z)[target])}. Zapisz
  ją w postaci stabilnej i powiedz, ile eksponent liczy każda wersja.
  \answerto{$\logsumexp(z) - z_{\text{target}}$. Wersja naiwna bierze
  eksponentę każdego wyniku, a potem logarytm ilorazu; wersja stabilna bierze
  eksponentę każdego przesuniętego wyniku raz, wewnątrz $\logsumexp$, a potem
  wykonuje jedno odejmowanie. Liczba eksponent jest ta sama, a problem z
  zakresem nie.}
\item Kompensacja Kahana opiera się na tym, że $(t - s) - y$ odzyskuje
  zaokrąglenie. Powiedz, jaka własność arytmetyki to umożliwia i kiedy zawodzi.
  \answerto{Dla dodawania zmiennoprzecinkowego błąd zaokrąglenia sumy $s + y$
  jest sam reprezentowalny dokładnie, gdy $|s| \ge |y|$, więc odejmowanie
  odzyskuje go bez własnego błędu. Zawodzi, gdy kompilatorowi wolno zmienić
  łączność wyrażenia, i dlatego ten algorytm zapisuje się często z barierami.}
\item Milion gradientów \code{fp16} sumuje się do akumulatora \code{fp16}.
  Oszacuj, jaka część przetrwa, i powiedz, jaka jedna zmiana naprawia większość.
  \answerto{Cokolwiek poniżej mniej więcej $\val{p01.swamp.fp16}$ bieżącej sumy
  nie wnosi nic, więc gdy suma urośnie, późniejsze składniki przepadają
  hurtem. Akumulowanie w \code{float32} naprawia prawie wszystko za cenę
  akumulatora, i dlatego jest to standardowy przepis.}
\item $\eps$ równe $10^{-8}$ dodaje się pod pierwiastkiem w normalizacji
  warstwowej działającej w \code{fp16}. Powiedz, co to osiąga.
  \answerto{Nic. Epsilon \code{fp16} jest rzędu $10^{-3}$, więc przy każdej
  wariancji, którą format reprezentuje, dodanie $10^{-8}$ nie zmienia żadnego
  bitu. $\eps$ trzeba przeskalować do formatu albo policzyć statystyki w
  szerszym.}
\end{furtherproblems}
